% Bagian: computability
% Bab: recursive-functions
% Subbagian: composition

\documentclass[../../../include/open-logic-section]{subfiles}

\begin{document}

\olfileid[id]{cmp}{rec}{com}
\olsection{Komposisi}

Jika $f$ dan $g$ adalah dua fungsi bilangan asli berargumen satu, kita dapat
mengomposisikannya: $h(x) = f(g(x))$. Fungsi baru~$h(x)$ kemudian
didefinisikan dengan \emph{komposisi} dari fungsi $f$ dan~$g$. Kita ingin
memperumum hal ini untuk fungsi dengan lebih dari satu argumen.

Berikut salah satu caranya: andaikan $f$ merupakan fungsi berargumen~$k$ dan
$g_0$, \dots, $g_{k-1}$ merupakan $k$ fungsi yang semuanya berargumen~$n$.
Kita dapat mendefinisikan suatu fungsi baru berargumen~$n$, yaitu $h$, sebagai
berikut:
\[
h(x_0, \dots, x_{n-1}) =
f(g_0(x_0, \dots, x_{n-1}), \dots, g_{k-1}(x_0, \dots, x_{n-1}))
\]
Jika $f$ dan semua $g_i$ dapat dihitung, maka $h$ juga dapat dihitung. Untuk
menghitung $h(x_0, \dots, x_{n-1})$, mula-mula hitung nilai
$y_i = g_i(x_0, \dots, x_{n-1})$ untuk setiap $i = 0$, \dots,~$k-1$.
Kemudian masukkan nilai-nilai tersebut ke dalam $f$ untuk menghitung
$h(x_0, \dots, x_{n-1}) = f(y_0, \dots, y_{k-1})$.

Pencirian ini mungkin tampak terlalu membatasi apa yang terjadi ketika kita
menghitung fungsi baru dengan menggunakan beberapa fungsi yang sudah ada.
Salah satu alasannya ialah kadang kita tidak menggunakan semua argumen suatu
fungsi, seperti ketika kita mendefinisikan $g(x, y, z) = \Succ(z)$ untuk
digunakan dalam definisi rekursif primitif bagi~$\Add$. Andaikan kita
diperbolehkan menggunakan fungsi-fungsi berikut:
\[
\Proj{n}{i}(x_0, \dots, x_{n-1}) = x_i
\]
Fungsi~$\Proj{n}{i}$ disebut fungsi \emph{proyeksi}: $\Proj{n}{i}$ merupakan
fungsi berargumen~$n$. Dengan demikian, $g$ dapat didefinisikan oleh
\[
g(x, y, z) = \Succ(\Proj{3}{2}(x, y, z)).
\]
Di sini, peran $f$ dimainkan oleh fungsi $1$-tempat~$\Succ$, sehingga
$k=1$. Kita mempunyai satu fungsi $3$-tempat $\Proj{3}{2}$ yang memainkan
peran $g_0$. Hasilnya adalah fungsi $3$-tempat yang mengembalikan penerus
argumen ketiga.

Fungsi proyeksi juga memungkinkan kita mendefinisikan fungsi baru dengan
mengurutkan kembali atau menyamakan argumen. Sebagai contoh, fungsi $h(x) =
\Add(x, x)$ dapat didefinisikan oleh
\[
h(x_0) = \Add(\Proj{1}{0}(x_0),\Proj{1}{0}(x_0)).
\]
Di sini $k=2$, $n=1$, peran $f(y_0,y_1)$ dimainkan oleh $\Add$, dan peran
$g_0(x_0)$ maupun $g_1(x_0)$ dimainkan oleh~$\Proj{1}{0}(x_0)$, yakni fungsi
proyeksi berargumen satu (atau fungsi identitas).

Jika $f(y_0, y_1)$ merupakan fungsi yang sudah kita miliki, kita dapat
mendefinisikan fungsi $h(x_0, x_1) = f(x_1, x_0)$ dengan
\[
h(x_0, x_1) = f(\Proj{2}{1}(x_0, x_1),\Proj{2}{0}(x_0, x_1)).
\]
Di sini $k=2$, $n = 2$, dan peran $g_0$ serta $g_1$ masing-masing dimainkan
oleh $\Proj{2}{1}$ dan~$\Proj{2}{0}$.

Anda mungkin juga mengkhawatirkan keharusan agar $g_0$, \dots,~$g_{k-1}$
semuanya mempunyai aritas yang sama, yaitu~$n$. (Ingat bahwa \emph{aritas}
suatu fungsi adalah banyaknya argumen; fungsi berargumen~$n$ mempunyai
aritas~$n$.) Namun, penambahan fungsi proyeksi memberikan keluwesan yang
diinginkan. Sebagai contoh, andaikan $f$ dan~$g$ merupakan fungsi berargumen~$3$,
sedangkan $h$ merupakan fungsi berargumen~$2$ yang didefinisikan oleh
\[
h(x,y) = f(x,g(x,x,y),y).
\]
Definisi~$h$ dapat ditulis ulang dengan fungsi proyeksi sebagai
\[
h(x,y) = f(\Proj{2}{0}(x,y), g(\Proj{2}{0}(x,y), \Proj{2}{0}(x,y),
\Proj{2}{1}(x,y)), \Proj{2}{1}(x,y)).
\]
Dengan demikian, $h$ merupakan komposisi $f$ dengan $\Proj{2}{0}$, $l$, dan
$\Proj{2}{1}$, dengan
\[
l(x,y) = g(\Proj{2}{0}(x,y),\Proj{2}{0}(x,y),\Proj{2}{1}(x,y)),
\]
yakni $l$ merupakan komposisi $g$ dengan $\Proj{2}{0}$, $\Proj{2}{0}$,
dan~$\Proj{2}{1}$.

\end{document}
